% GENERATED FILE. Edit canonical lessons, then run npm run generate:books.
% canonical-source-hash: fnv1a64:64a7e7cee72262b2
% canonical-lessons: GE-C339-aeusserung, GE-C339-hinweisen, GE-C339-hinweis, GE-C339-hervorheben, GE-C339-darlegen

\chapter{Etwas hervorheben}
\label{ch:ge-etwas-hervorheben}
\hlchaptermodality{339}

\begin{chapteropening}
\textbf{By the end of this chapter:} \emph{I can make a remark, call attention to something, single it out and set out my reasons.}

Complete the last lesson of chapter 339: I can make a remark, call attention to something, single it out and set out my reasons.
\end{chapteropening}

\section[die Äußerung]{die Äußerung — the remark, the statement}

\label{lesson:GE-C339-aeusserung}



\begin{quote}
Before the new one: say the German for obviously, then the German for to suspect.
\end{quote}

\subsection*{You'll want to know: die Äußerung}
\textbf{die Äußerung} — \textquotedblleft{}the remark, the statement\textquotedblright{}.

\begin{quote}
Seine Äußerung über die Kollegen war nicht nett.
\end{quote}

\begin{quote}
\emph{His remark about his colleagues was not nice.}
\end{quote}

Plural: \emph{die Äußerungen}. The noun of \emph{äußern} — to express.

\begin{cousinweb}
\emph{Äußerung} is the noun of \emph{äußern} \textquotedblleft{}to express\textquotedblright{}, built on \emph{außer}, from Middle High German \emph{ūzer} and Old High German \emph{ūzar} \textquotedblleft{}out of\textquotedblright{} (EtymDB): you put a thought outside.
\end{cousinweb}

\begin{grammarlens}[title={What you've built}]
One more word for telling, explaining and describing at B1.
\end{grammarlens}

\subsection*{Your turn}
\begin{itemize}
\raggedright
  \item \emph{Say it:} \emph{die Äußerung}
  \item \emph{Say it:} \emph{die Äußerung}, once more
  \item \emph{Say it:} a sentence of your own with \emph{die Äußerung}
  \item \emph{Recall:} say the German for obviously, then the German for to suspect, then say \emph{die Äußerung} again
\end{itemize}

\begin{tcolorbox}[breakable,colback=teal!4,colframe=teal!35!black,title={Before you move on}]
What does die Äußerung mean? (\textquotedblleft{}The remark, the statement\textquotedblright{}.) Say it once more.
\end{tcolorbox}
\section[hinweisen]{hinweisen — to indicate, to call attention to}

\label{lesson:GE-C339-hinweisen}



\begin{quote}
Before the new one: say the German for to suspect, then the German for the remark.
\end{quote}

\subsection*{You'll want to know: hinweisen}
\textbf{hinweisen} — \textquotedblleft{}to indicate, to call attention to\textquotedblright{}.

\begin{quote}
Der Arzt hat darauf hingewiesen, dass ich mehr Wasser trinken soll.
\end{quote}

\begin{quote}
\emph{The doctor noted that I should drink more water.}
\end{quote}

Separable and strong: \emph{er weist darauf hin}, \emph{er hat darauf hingewiesen}. \emph{Auf etwas hinweisen} — to call attention to something.

\begin{cousinweb}
\emph{hinweisen} is built from \emph{hin} \textquotedblleft{}towards there\textquotedblright{} and \emph{weisen} \textquotedblleft{}to show, to direct\textquotedblright{} (EtymDB): you show someone where to look.
\end{cousinweb}

\begin{grammarlens}[title={What you've built}]
One more word for telling, explaining and describing at B1.
\end{grammarlens}

\subsection*{Your turn}
\begin{itemize}
\raggedright
  \item \emph{Say it:} \emph{hinweisen}
  \item \emph{Say it:} \emph{hinweisen}, once more
  \item \emph{Say it:} a sentence of your own with \emph{hinweisen}
  \item \emph{Recall:} say the German for to suspect, then the German for the remark, then say \emph{hinweisen} again
\end{itemize}

\begin{tcolorbox}[breakable,colback=teal!4,colframe=teal!35!black,title={Before you move on}]
What does hinweisen mean? (\textquotedblleft{}To indicate, to call attention to\textquotedblright{}.) Say it once more.
\end{tcolorbox}
\section[der Hinweis]{der Hinweis — the hint, the tip}

\label{lesson:GE-C339-hinweis}



\begin{quote}
Before the new one: say the German for the remark, then the German for to indicate.
\end{quote}

\subsection*{You'll want to know: der Hinweis}
\textbf{der Hinweis} — \textquotedblleft{}the hint, the tip\textquotedblright{}.

\begin{quote}
Danke für den Hinweis! Ich hatte den Fehler nicht bemerkt.
\end{quote}

\begin{quote}
\emph{Thanks for the tip! I had not noticed the mistake.}
\end{quote}

Plural: \emph{die Hinweise}. On a sign, \emph{Hinweis} is a notice.

\begin{cousinweb}
\emph{Hinweis} is the noun of \emph{hinweisen}, built from \emph{hin} \textquotedblleft{}towards there\textquotedblright{} and \emph{weisen} \textquotedblleft{}to show, to direct\textquotedblright{} (EtymDB).
\end{cousinweb}

\begin{grammarlens}[title={What you've built}]
One more word for telling, explaining and describing at B1.
\end{grammarlens}

\subsection*{Your turn}
\begin{itemize}
\raggedright
  \item \emph{Say it:} \emph{der Hinweis}
  \item \emph{Say it:} \emph{der Hinweis}, once more
  \item \emph{Say it:} a sentence of your own with \emph{der Hinweis}
  \item \emph{Recall:} say the German for the remark, then the German for to indicate, then say \emph{der Hinweis} again
\end{itemize}

\begin{tcolorbox}[breakable,colback=teal!4,colframe=teal!35!black,title={Before you move on}]
What does der Hinweis mean? (\textquotedblleft{}The hint, the tip\textquotedblright{}.) Say it once more.
\end{tcolorbox}
\section[hervorheben]{hervorheben — to bring out, to single out}

\label{lesson:GE-C339-hervorheben}



\begin{quote}
Before the new one: say the German for to indicate, then the German for the hint.
\end{quote}

\subsection*{You'll want to know: hervorheben}
\textbf{hervorheben} — \textquotedblleft{}to bring out, to single out\textquotedblright{}.

\begin{quote}
Im Bericht hebt die Lehrerin hervor, dass meine Tochter sehr fleißig ist.
\end{quote}

\begin{quote}
\emph{In the report the teacher singles out how hard-working my daughter is.}
\end{quote}

Separable and strong: \emph{sie hebt hervor}, \emph{sie hat hervorgehoben}. Literally to lift something out to the front.

\begin{cousinweb}
\emph{hervorheben} is built from \emph{hervor} \textquotedblleft{}forth, out\textquotedblright{} and \emph{heben} \textquotedblleft{}to lift\textquotedblright{}, inherited from Old High German \emph{heffen} (EtymDB): you lift a point out so that everyone notices it.
\end{cousinweb}

\begin{grammarlens}[title={What you've built}]
One more word for telling, explaining and describing at B1.
\end{grammarlens}

\subsection*{Your turn}
\begin{itemize}
\raggedright
  \item \emph{Say it:} \emph{hervorheben}
  \item \emph{Say it:} \emph{hervorheben}, once more
  \item \emph{Say it:} a sentence of your own with \emph{hervorheben}
  \item \emph{Recall:} say the German for to indicate, then the German for the hint, then say \emph{hervorheben} again
\end{itemize}

\begin{tcolorbox}[breakable,colback=teal!4,colframe=teal!35!black,title={Before you move on}]
What does hervorheben mean? (\textquotedblleft{}To bring out, to single out\textquotedblright{}.) Say it once more.
\end{tcolorbox}
\section[darlegen]{darlegen — to set out, to explain clearly}

\label{lesson:GE-C339-darlegen}



\begin{quote}
Before the new one: say the German for the hint, then the German for to bring out.
\end{quote}

\subsection*{You'll want to know: darlegen}
\textbf{darlegen} — \textquotedblleft{}to set out, to explain clearly\textquotedblright{}.

\begin{quote}
Sie legt ihre Gründe ruhig und klar dar.
\end{quote}

\begin{quote}
\emph{She sets out her reasons calmly and clearly.}
\end{quote}

Separable: \emph{sie legt dar}, \emph{sie hat dargelegt}. Formal, often in writing.

\begin{cousinweb}
\emph{darlegen} is built from \emph{dar} \textquotedblleft{}there, forth\textquotedblright{} and \emph{legen} \textquotedblleft{}to lay\textquotedblright{}, inherited from Old High German \emph{leggen} and Proto-Germanic \emph{lagjaną} (EtymDB): you lay your reasons out in front of people.
\end{cousinweb}

\begin{grammarlens}[title={What you've built}]
One more word for telling, explaining and describing at B1.
\end{grammarlens}

\subsection*{Your turn}
\begin{itemize}
\raggedright
  \item \emph{Say it:} \emph{darlegen}
  \item \emph{Say it:} \emph{darlegen}, once more
  \item \emph{Say it:} a sentence of your own with \emph{darlegen}
  \item \emph{Recall:} say the German for the hint, then the German for to bring out, then say \emph{darlegen} again
\end{itemize}

\begin{tcolorbox}[breakable,colback=teal!4,colframe=teal!35!black,title={Before you move on}]
What does darlegen mean? (\textquotedblleft{}To set out, to explain clearly\textquotedblright{}.) Say it once more.
\end{tcolorbox}
